\section{GitHub Connection \& Repository Management}

\subsection{Introduction}

GitHub is used as the source code management platform for the LivingDocs project. All documentation, source code, and change history are centrally stored on GitHub to support team collaboration more effectively.

In the Software Engineering project, integrating GitHub with the \textbf{LivingDocs} system is necessary to achieve the following goals:

\begin{itemize}
    \item Centralized and consistent source code management among team members.
    \item Synchronize source code changes with the automated documentation system.
    \item Track commit history and Pull Requests to control development progress.
    \item Support review processes and software quality checks before deployment.
    \item Reduce outdated documentation compared to actual source code.
\end{itemize}

Thanks to direct integration with GitHub, LivingDocs can automatically detect source code changes and update corresponding documentation, maintaining synchronization between technical documents and the evolving system.

\subsection{GitHub Connection Using OAuth2}

To allow users to connect their GitHub accounts to LivingDocs without providing passwords directly, the system uses the OAuth2 Authorization Code Flow supported by GitHub.

\paragraph{OAuth2 Mechanism}

OAuth2 allows users to grant LivingDocs access without sharing their GitHub password. The system only receives an \texttt{access token} with the scope approved by the user.

\paragraph{Login and Authentication Process}

\begin{enumerate}
    \item The user selects the \textbf{Connect GitHub} function.
    \item The system redirects to GitHub’s authentication page.
    \item The user logs in to GitHub and grants repository access.
    \item GitHub returns an \texttt{authorization code}.
    \item LivingDocs sends this code to GitHub to obtain an \texttt{access token}.
    \item The token is encrypted and securely stored in the system.
\end{enumerate}

\paragraph{Benefits of OAuth2}

\begin{itemize}
    \item \textbf{High security}: GitHub passwords are not stored in the system.
    \item \textbf{Clear access control}: Only authorized repositories are accessible.
    \item \textbf{Easy revocation}: Users can revoke access at any time on GitHub.
    \item \textbf{Modern compliance}: Compatible with GitHub’s security policies.
\end{itemize}

This mechanism ensures LivingDocs maintains strong GitHub integration while keeping user data and accounts secure.

\subsection{Repository Management}

After successfully connecting to GitHub, LivingDocs allows users to select repositories to manage and synchronize with the documentation system.

\paragraph{Repository Selection}

The system uses the GitHub REST API to retrieve the list of repositories the user can access. Users can select one or multiple repositories to include in the LivingDocs workspace.

Stored information includes:

\begin{itemize}
    \item Repository ID.
    \item Repository name.
    \item Owner.
    \item Default branch.
    \item Last synchronization time.
\end{itemize}

\paragraph{Data Synchronization from Main Branch}

Once connected, users can choose repositories to manage. Repository information is stored for tracking and synchronization:

\begin{itemize}
    \item Update the latest commit.
    \item Retrieve the list of changed files.
    \item Synchronize merged Pull Requests.
    \item Update repository metadata.
\end{itemize}

\paragraph{Update Status Tracking}

The system maintains synchronization status through indicators such as:

\begin{itemize}
    \item Last processed commit.
    \item Open Pull Requests.
    \item Merged Pull Requests.
    \item Last synchronization time.
    \item Success or failure status of synchronization.
\end{itemize}

Currently, this content is at the design and solution description stage. Implementation of repository connection and synchronization will be carried out once the Backend and Database components are completed.

\subsection{CI/CD Integration with GitHub Actions}

To support development in later stages, LivingDocs plans to integrate with \textbf{GitHub Actions} to automate the \textbf{Continuous Integration / Continuous Deployment (CI/CD)} process. This integration will be implemented once Backend, Frontend, and Database components are ready.

\paragraph{Automated Build/Test Workflow}

During system deployment, when a new commit is pushed or a Pull Request is created, GitHub Actions can be configured to trigger automated workflows. The process typically includes:

\begin{enumerate}
    \item Checkout source code.
    \item Set up runtime environment.
    \item Build the project.
    \item Run unit tests and integration tests.
    \item Send results back to GitHub.
\end{enumerate}

\paragraph{Automatic Documentation Update on PR Merge}

After a Pull Request is merged into the main branch, the workflow can call LivingDocs API to:

\begin{itemize}
    \item Analyze source code changes.
    \item Identify related documentation to update.
    \item Generate documentation update proposals.
    \item Synchronize new documentation versions into the system.
\end{itemize}

\paragraph{Benefits of CI/CD}

\begin{itemize}
    \item \textbf{Reduce manual errors}: Avoid forgetting documentation updates or missing checks.
    \item \textbf{Accelerate development}: Automated checks detect errors earlier.
    \item \textbf{Ensure quality}: All changes are built and tested before merging.
    \item \textbf{Maintain synchronization}: Documentation is always updated with source code changes.
\end{itemize}

In later stages, the team plans to integrate GitHub Actions to support CI/CD workflows.

\subsection{Pull Request and Commit Tracking}

Tracking Pull Requests and commits is essential for LivingDocs to link source code changes with corresponding technical documentation.

\paragraph{Recording Changes from Pull Requests}

When GitHub sends a \texttt{pull\_request} event, the system will:

\begin{enumerate}
    \item Receive the webhook from GitHub.
    \item Verify the webhook’s security signature.
    \item Retrieve the list of added, modified, or deleted files.
    \item Classify the impact level on documentation.
    \item Create a tracking record in LivingDocs.
\end{enumerate}

\paragraph{Linking Commits to Documentation}

Each commit is associated with:

\begin{itemize}
    \item Affected system modules or components.
    \item Changed APIs or classes.
    \item Related design or API documentation.
    \item Documentation version requiring update.
\end{itemize}

This mechanism enables traceability between documentation and commits, supporting auditing and system maintenance.

\paragraph{Review and Merge Workflow}

\begin{enumerate}
    \item Developer creates a Pull Request.
    \item GitHub Actions automatically build and run tests.
    \item LivingDocs analyzes documentation impact.
    \item Reviewer checks source code and proposes documentation updates.
    \item Once approved, the Pull Request is merged into the main branch.
    \item The system triggers final documentation synchronization.
\end{enumerate}

This workflow ensures all important source code changes are reviewed alongside corresponding documentation before integration.

\subsection{Security and Error Handling}

Since GitHub Integration works with external services and critical source code, the system requires proper security and error handling mechanisms to ensure stability and safety.

\paragraph{Expired Tokens}

GitHub access tokens may expire or be revoked. When detected, LivingDocs will:

\begin{itemize}
    \item Mark connection status as \texttt{EXPIRED}.
    \item Pause repository synchronization.
    \item Display a message requesting the user to reconnect GitHub.
\end{itemize}

\paragraph{GitHub Connection Errors}

Common errors include:

\begin{itemize}
    \item Network disconnection.
    \item GitHub API not responding.
    \item Exceeding request limits (rate limit).
    \item Webhook delivery failure.
\end{itemize}

If GitHub API does not respond, the system retries multiple times before notifying the user.

\paragraph{Warning and Token Refresh Mechanism}

LivingDocs supports:

\begin{itemize}
    \item Alerts when tokens are about to expire.
    \item Notifications to users via the system interface.
    \item If tokens are invalid or revoked, the system requests reconnection to GitHub.
    \item Logging all authentication and synchronization processes for troubleshooting.
\end{itemize}

\paragraph{Data Security During Synchronization}

To protect source code and credentials, the system applies:

\begin{itemize}
    \item Tokens stored in encrypted form to prevent leaks.
    \item HTTPS/TLS for all GitHub connections.
    \item Verification of webhook \texttt{X-Hub-Signature-256} to prevent forgery.
    \item Role-Based Access Control (RBAC).
    \item No logging of tokens or sensitive data in application logs.
\end{itemize}

This chapter provides the foundation for implementing GitHub Integration in upcoming sprints. At the current stage, the team focuses on requirements, system design, and technical documentation before starting source code development.
